\documentclass[11pt]{article}
\usepackage{amsmath,amsthm,amssymb}
\usepackage[margin=1in]{geometry}
\usepackage{enumitem}
\usepackage{booktabs}

\newtheorem{theorem}{Theorem}
\newtheorem{lemma}[theorem]{Lemma}
\newtheorem{proposition}[theorem]{Proposition}
\newtheorem{corollary}[theorem]{Corollary}
\newtheorem{conjecture}[theorem]{Conjecture}
\theoremstyle{definition}
\newtheorem{definition}[theorem]{Definition}
\newtheorem{remark}[theorem]{Remark}
\newtheorem{example}[theorem]{Example}

\title{The rank-zero theorem for $\lambda=(3,2,1^{n-5})$ \\
  via a three-branch reduction}
\author{Clio}
\date{20 May 2026}

\begin{document}
\maketitle

\begin{abstract}
We prove $\Pi^{S_n}\!\mid_{V_{(3,2,1^{n-5})}}\!=0$ for every
$n\ge 8$.  This is the first family in the rank-zero programme
where $\lambda$ has \emph{two} length-preserving corners, so
$B_{\ell+1}^{(\lambda)} V_\lambda$ is $2$-dimensional rather than
$1$-dimensional.  Empirically the $j$-formula at $j=\ell+1$ does
not collapse the image to dimension $1$, which was the obstacle
flagged in the May-19 note.  We resolve it by showing the
$2$-dimensional image is nevertheless \emph{entirely contained in
$E_1^-$}, so the rightmost $(T_1{+}1)$ of $R'_{\ell}$ annihilates
it.

The proof is a three-branch generalization of the May-15
branching reduction. The +q-eigenspace
$E_{n-1}^+\!\mid_{V_\lambda}$ decomposes, as an
$H_q(S_{n-2})$-module, into \emph{three} pieces
$\Phi_A(V_{\mu_A})\oplus\Phi_B(V_{\mu_B})\oplus\Phi_C(V_{\mu_C})$
indexed by the three pairs of corners of $\lambda$.  Two pieces
(those involving the column-$1$ bottom corner) reduce to the
known $j$-formula for $\mu_A=(2,2,1^{n-6})$ (May-15/18 family) and
$\mu_B=(3,1^{n-5})$ (May-12 hook), each contributing a
$1$-dimensional image in $E_1^-$.  The third piece, indexed by
$\mu_C=(2,1^{n-4})$ (May-11 hook), contributes \emph{zero} because
its $j$-formula sits one step earlier, so the May-15 commutation
crosses the $E_1^-$ wall and is killed by a leading $(T_1{+}1)$.

\textbf{Status.}  Conditional only on Phase A endpoint at $\lambda$
(immediate from May-19 abstract Phase A theorem).  All other
ingredients are rigorously proved in prior writeups.

\textbf{Verification.}  $B_{\ell+1}^{(\lambda)} V_\lambda\subseteq
E_1^-$ verified at $q=7$ for $n\in\{8,9,10\}$ (dim $2$ each, full
collapse to $0$ after one more $R'_\ell$). Dimension decomposition
$\dim V_{\mu_A}+\dim V_{\mu_B}+\dim V_{\mu_C}=\dim E_{n-1}^+ V_\lambda$
verified through $n=11$.
\end{abstract}

\section{Setup}\label{sec:setup}

We use the conventions of the May-13--May-19 papers.  $H_q(S_n)$
is the Hecke algebra over $\mathbb{Q}(q)$ with generators
$T_1,\ldots,T_{n-1}$, quadratic relation $T_i^2=(q-1)T_i+q$, and
braid relations.  $V_\lambda$ is the seminormal cell module for
$\lambda\vdash n$, with Hoefsmit basis $\{v_T\}_{T\in\mathrm{SYT}(\lambda)}$.
$E_j^+$ (resp.\ $E_j^-$) is the $q$-eigenspace (resp.\ $-1$-eigenspace)
of $T_j$.

Write
\[
   R'_k=(T_{k-1}{+}1)(T_{k-2}{+}1)\cdots(T_1{+}1),
   \qquad
   B_j^{(\lambda)}=R'_j R'_{j+1}\cdots R'_n,
   \qquad
   \Pi^{S_n}=B_2^{(\lambda)}.
\]

Throughout, $n\ge 8$ and
\[
   \lambda:=(3,2,1^{n-5})\vdash n,
   \qquad \ell:=\ell(\lambda)=n-3,
   \qquad j_0:=\ell+1=n-2.
\]
The $j$-formula target is
$B_{n-2}^{(\lambda)}=R'_{n-2}R'_{n-1}R'_n$ ($3$ factors).

\subsection{The three corners and the three inner partitions}

$\lambda=(3,2,1^{n-5})$ has \emph{three} removable corners:
\[
   c_1=(1,3),\qquad c_2=(2,2),\qquad c_3=(n{-}3,1).
\]
Of these, $c_1$ and $c_2$ are \emph{length-preserving}
(removing them leaves a partition of length $\ell-1$ but with the
same number of rows occupied as $\lambda$ ignoring trailing parts;
more precisely, $\lambda_r>\lambda_{r+1}$ at $r=1,2$), and $c_3$ is
the column-$1$ bottom corner (length-reducing).  The rank-corner
formula (May-11) predicts $\dim B_{\ell+1}^{(\lambda)} V_\lambda=
\#\{\text{length-pres corners}\}=2$; this matches our computation
(see Section~\ref{sec:verify}).

For each \emph{pair} of corners $\{c_i,c_j\}$, the doubly-stripped
partition (remove both cells) is
\begin{align*}
   \mu_A &:= \lambda\setminus\{c_1,c_3\}=(2,2,1^{n-6})\vdash n-2,
   &\text{(``A''=length-pres $c_1$ paired with bottom $c_3$)} \\
   \mu_B &:= \lambda\setminus\{c_2,c_3\}=(3,1^{n-5})\vdash n-2,
   &\text{(``B''=length-pres $c_2$ paired with bottom $c_3$)} \\
   \mu_C &:= \lambda\setminus\{c_1,c_2\}=(2,1^{n-4})\vdash n-2.
   &\text{(``C''=the two length-pres corners with each other)}
\end{align*}
$\mu_A$ is the second non-hook (May-15 family at $a=2$), $\mu_B$ is
the column-$1$ hook of length $\ge 3$ in the first row (May-12),
$\mu_C$ is the column-$1$ hook of length $2$ in the first row
(May-11).

\subsection{Known rank-zero results we will use}

Each $\mu_X$ is in its own rank-zero regime, and the rank-zero
theorem is known structurally:

\begin{theorem}[$j$-formula at $\mu_A$, May-13 + May-18]\label{thm:muA}
For $\mu_A=(2,2,1^{m-4})\vdash m$ with $m\ge 6$,
$B_{m-1}^{(\mu_A)} V_{\mu_A}$ is $1$-dimensional and contained in
$E_1^-\!\mid_{V_{\mu_A}}$.
\end{theorem}

\begin{theorem}[$j$-formula at $\mu_B$, May-12 hook]\label{thm:muB}
For $\mu_B=(3,1^{m-3})\vdash m$ with $m\ge 6$,
$B_{m-1}^{(\mu_B)} V_{\mu_B}$ is $1$-dimensional and contained in
$E_1^-\!\mid_{V_{\mu_B}}$.
\end{theorem}

\begin{theorem}[$j$-formula at $\mu_C$, May-11 hook]\label{thm:muC}
For $\mu_C=(2,1^{m-2})\vdash m$ with $m\ge 4$,
$B_m^{(\mu_C)} V_{\mu_C}=R'^{(\mu_C)}_m V_{\mu_C}$ is
$1$-dimensional and contained in $E_1^-\!\mid_{V_{\mu_C}}$.
\end{theorem}

\begin{theorem}[Abstract Phase A endpoint, May-19]\label{thm:phaseA}
For any $H_q(S_n)$-module $V$ and any $1\le j\le n-1$,
$W_j(V):=(T_j{+}1)\cdots(T_1{+}1)V=E_j^+(V)$.  In particular,
$R'_n V=E_{n-1}^+(V)$.
\end{theorem}

We apply Theorem~\ref{thm:phaseA} at $V=V_\lambda$ and
at $V=V_{\mu_X}$ ($X\in\{A,B,C\}$); we apply
Theorems~\ref{thm:muA}--\ref{thm:muC} at $m=n-2$, noting that
$m\ge 6$ for $n\ge 8$.

\section{The branching isomorphisms $\Phi_X$}\label{sec:Phi}

We mimic the May-15 construction for each pair of corners.  Recall
that $H_q(S_{n-2})$ embeds in $H_q(S_n)$ as the subalgebra
generated by $\{T_1,\ldots,T_{n-3}\}$, and $T_{n-1}$ commutes with
every $T_i$ for $i\le n-3$ (index gap $\ge 2$).  Hence
$E_{n-1}^+\!\mid_{V_\lambda}$ is an $H_q(S_{n-2})$-submodule of
$V_\lambda$.

\begin{definition}[$\Phi_X$, $X\in\{A,B,C\}$]\label{def:Phi}
Fix $X\in\{A,B,C\}$ and write $\{c_X^-,c_X^+\}=\{c_i,c_j\}$ as the
ordered pair (by row, lower-row first) corresponding to $\mu_X$:
\[
   X=A:\ (c_A^-,c_A^+)=(c_1,c_3)=((1,3),(n-3,1));\qquad
   X=B:\ (c_2,c_3)=((2,2),(n-3,1));
\]
\[
   X=C:\ (c_1,c_2)=((1,3),(2,2)).
\]
For each SYT $T'$ of $\mu_X$ on letters $\{1,\ldots,n-2\}$, define
two SYTs $T^X_A,T^X_B$ of $\lambda$ on letters $\{1,\ldots,n\}$ by
\begin{itemize}[leftmargin=2em,topsep=0.3em,itemsep=0.2em]
\item $T^X_A$: extend $T'$ by placing $n-1$ at $c_X^-$ and $n$ at
$c_X^+$.
\item $T^X_B$: extend $T'$ by placing $n$ at $c_X^-$ and $n-1$ at
$c_X^+$.
\end{itemize}
Both are valid SYTs of $\lambda$ (since $n-1,n$ exceed all other
entries, the column- and row-increase conditions involving the
added cells are automatic; the cells $c_X^\pm$ are corners of
$\lambda$ once $T'$ has been placed).

The axial distance $d_X:=c(c_X^+)-c(c_X^-)$ (content difference) is
non-zero of absolute value $\ge 2$ in all three cases:
\[
   d_A=(1-(n-3))-(3-1)=4-n,\quad
   d_B=(1-(n-3))-(2-2)=4-n,\quad
   d_C=(2-2)-(3-1)=-2.
\]
For $n\ge 8$, $|d_A|=|d_B|=n-4\ge 4$ and $|d_C|=2$, so all
$T_{n-1}$-$2$-blocks $\{v_{T^X_A},v_{T^X_B}\}$ are non-degenerate
over $\mathbb{Q}(q)$ (the Hoefsmit ratio $\tau(d)\in\mathbb{Q}(q)^\times$
for $|d|\ge 2$).

Define
\[
   \Phi_X(v_{T'}^{(\mu_X)})\;:=\;
   \text{the $T_{n-1}$ $+q$-eigenvector in the $2$-block}\
   \{v_{T^X_A},v_{T^X_B}\}\subset V_\lambda,
\]
normalized so the coefficient of $v_{T^X_A}$ equals $1$.  Extend
linearly to $\Phi_X:V_{\mu_X}\to V_\lambda$.
\end{definition}

\begin{lemma}[Equivariance of $\Phi_X$]\label{lem:Phi-equiv}
For every $X\in\{A,B,C\}$, every $1\le i\le n-3$ and every
$v\in V_{\mu_X}$,
$T_i\,\Phi_X(v)=\Phi_X(T_i\,v)$.
\end{lemma}

\begin{proof}
By linearity, suffices to check $v=v_{T'}^{(\mu_X)}$.  The
positions of $i,i{+}1$ in $T^X_A$ and $T^X_B$ are identical to the
positions in $T'$ (since $i,i{+}1\le n-2$ and the cells of
$T^X_A,T^X_B$ holding entries $\le n-2$ agree with $T'$).  Three
cases for the $T_i$-action on $T'$ (the argument is verbatim
May-15 Lemma~3.2 in our notation):

\noindent\textbf{(a) Same row in $T'$.}  $T_iv_{T'}^{(\mu_X)}=
qv_{T'}^{(\mu_X)}$.  Since the positions of $i,i{+}1$ in
$T^X_A,T^X_B$ are the same as in $T'$, both $T_iv_{T^X_A}=qv_{T^X_A}$
and $T_iv_{T^X_B}=qv_{T^X_B}$, hence $T_i\Phi_X(v_{T'}^{(\mu_X)})=
q\Phi_X(v_{T'}^{(\mu_X)})$.

\noindent\textbf{(b) Same column in $T'$.}  Symmetric to (a),
eigenvalue $-1$.

\noindent\textbf{(c) $T_i$ $2$-block on $V_{\mu_X}$.}  The block
partner is $s_iT'$ with axial distance $d_i$.  $T_i$ commutes with
$T_{n-1}$, so the $4$-dim subspace
$\mathrm{span}(v_{T^X_A},v_{T^X_B},v_{s_iT^X_A},v_{s_iT^X_B})$ is
preserved by both, and the actions decouple:
\begin{itemize}[leftmargin=2em,topsep=0.3em,itemsep=0.2em]
\item $T_{n-1}$: $2$-blocks $\{v_{T^X_A},v_{T^X_B}\}$ and
$\{v_{s_iT^X_A},v_{s_iT^X_B}\}$, axial distance $d_X$ (same in both).
\item $T_i$: $2$-blocks $\{v_{T^X_A},v_{s_iT^X_A}\}$ and
$\{v_{T^X_B},v_{s_iT^X_B}\}$, axial distance $d_i$ (same in both).
\end{itemize}
The Hoefsmit coefficients on each $T_{n-1}$-$2$-block are
determined by $d_X$ alone, so the same $\tau_X:=\tau(d_X)$ appears
in both block-partner relations.  Similarly the same $a_i,b_i$
appear in $T_iv_{T^X_A}=a_iv_{T^X_A}+b_iv_{s_iT^X_A}$ and
$T_iv_{T^X_B}=a_iv_{T^X_B}+b_iv_{s_iT^X_B}$ (because the axial
distance $d_i$ is the same).  Hence:
\begin{align*}
   T_i\Phi_X(v_{T'}^{(\mu_X)})
   &= T_i(v_{T^X_A}+\tau_X v_{T^X_B})
   = a_iv_{T^X_A}+b_iv_{s_iT^X_A}+\tau_X(a_iv_{T^X_B}+b_iv_{s_iT^X_B})\\
   &= a_i(v_{T^X_A}+\tau_Xv_{T^X_B})+b_i(v_{s_iT^X_A}+\tau_Xv_{s_iT^X_B})\\
   &= a_i\,\Phi_X(v_{T'}^{(\mu_X)})+b_i\,\Phi_X(v_{s_iT'}^{(\mu_X)})
   = \Phi_X(T_iv_{T'}^{(\mu_X)}).\qedhere
\end{align*}
\end{proof}

\begin{lemma}[$\Phi_X$ is injective]\label{lem:Phi-inj}
For each $X\in\{A,B,C\}$, $\Phi_X:V_{\mu_X}\hookrightarrow V_\lambda$
is injective.  Moreover, the seminormal supports of
$\Phi_A(V_{\mu_A}),\Phi_B(V_{\mu_B}),\Phi_C(V_{\mu_C})$ are pairwise
disjoint, so the sum
\[
   \Phi_A(V_{\mu_A})+\Phi_B(V_{\mu_B})+\Phi_C(V_{\mu_C})
\]
is direct.
\end{lemma}

\begin{proof}
For injectivity, observe that distinct $T',T''$ of $\mu_X$ give
disjoint pairs $\{T^X_A,T^X_B\}$ (the cells of $n-1,n$ are fixed at
$c_X^\pm$, so $\mu_X$ is recovered by deletion).  Hence
$\{\Phi_X(v_{T'}^{(\mu_X)})\}_{T'\in\mathrm{SYT}(\mu_X)}$ is
linearly independent in $V_\lambda$ (each is supported on a unique
disjoint $2$-element subset of the Hoefsmit basis of $V_\lambda$).

For pairwise disjointness across $X$:  $\Phi_X(v_{T'}^{(\mu_X)})$ is
supported on SYTs with $\{n-1,n\}$ at the cell pair
$\{c_X^-,c_X^+\}$.  Across $X\in\{A,B,C\}$, these cell pairs are
distinct, so the supports are pairwise disjoint.
\end{proof}

\begin{lemma}[Image equals $E_{n-1}^+\!\mid_{V_\lambda}$]\label{lem:Phi-image}
$\Phi_A(V_{\mu_A})\oplus\Phi_B(V_{\mu_B})\oplus\Phi_C(V_{\mu_C})
=E_{n-1}^+\!\mid_{V_\lambda}$.
\end{lemma}

\begin{proof}
\emph{Containment in $E_{n-1}^+$.}  By construction each
$\Phi_X(v_{T'}^{(\mu_X)})$ is a $T_{n-1}$ $+q$-eigenvector.

\emph{Equality by dimension count.}  Every $T_{n-1}$ $+q$-eigenvector
in the seminormal basis of $V_\lambda$ comes from a (necessarily
non-degenerate) $2$-block where $\{n-1,n\}$ occupy two different
corners of $\lambda$ at different rows and columns.  The three
corner pairs are $\{c_1,c_3\}=A,\ \{c_2,c_3\}=B,\ \{c_1,c_2\}=C$
(all in different rows/cols).  The number of SYTs whose
$\{n-1,n\}$ occupy a fixed corner pair $\{c^-,c^+\}$ equals
$2\cdot\#\mathrm{SYT}(\lambda\setminus\{c^-,c^+\})$, contributing
$\#\mathrm{SYT}(\mu_X)=\dim V_{\mu_X}$ to the $+q$-eigenspace
(one $+q$-eigenvector per $2$-block).  Summing,
\[
   \dim E_{n-1}^+\!\mid_{V_\lambda}
   =\dim V_{\mu_A}+\dim V_{\mu_B}+\dim V_{\mu_C}.
\]
By Lemma~\ref{lem:Phi-inj} the direct sum on the LHS of the
statement has the same dimension and is contained in
$E_{n-1}^+\!\mid_{V_\lambda}$, hence equals it.
\end{proof}

\begin{remark}\label{rem:dim-check}
At $n=8$: $\dim V_{\mu_A}|_6+\dim V_{\mu_B}|_6+\dim V_{\mu_C}|_6
=9+10+5=24=\dim E_{n-1}^+\!\mid_{V_\lambda}$ (computed at $q=7$
in Section~\ref{sec:verify}).  Verified through $n=11$
($27+28+8=63$).
\end{remark}

\begin{lemma}[$\Phi_X$ preserves $E_1^-$]\label{lem:Phi-E1}
For each $X$, $\Phi_X(E_1^-\!\mid_{V_{\mu_X}})\subseteq
E_1^-\!\mid_{V_\lambda}$.
\end{lemma}

\begin{proof}
$T_1$ is in $H_q(S_{n-2})$ and $\Phi_X$ is $T_1$-equivariant
(Lemma~\ref{lem:Phi-equiv}). Hence $T_1v=-v$ implies
$T_1\Phi_X(v)=-\Phi_X(v)$.
\end{proof}

\section{The reduction formula}\label{sec:reduction}

By Theorem~\ref{thm:phaseA} applied at $V=V_\lambda$,
$R'_n V_\lambda=E_{n-1}^+\!\mid_{V_\lambda}$, and by
Lemma~\ref{lem:Phi-image} this equals $\bigoplus_X\Phi_X(V_{\mu_X})$.

\begin{theorem}[Three-branch reduction]\label{thm:reduction}
\begin{equation}\label{eq:reduction}
   B_{n-2}^{(\lambda)} V_\lambda
   \;=\;
   (T_{n-3}{+}1)(T_{n-2}{+}1)\,
   \Big(\Phi_A(B_{n-3}^{(\mu_A)} V_{\mu_A})
       \;+\;\Phi_B(B_{n-3}^{(\mu_B)} V_{\mu_B})
       \;+\;\Phi_C(B_{n-3}^{(\mu_C)} V_{\mu_C})\Big).
\end{equation}
\end{theorem}

\begin{proof}
We follow the May-15 argument verbatim, applied to each
$\Phi_X$ piece simultaneously.

\paragraph{Step 1.}  $R'_n V_\lambda=\bigoplus_X\Phi_X(V_{\mu_X})$,
proved above.

\paragraph{Step 2.}  $R'_{n-1}=(T_{n-2}{+}1)R'_{n-2}^{\,[n-1]}$
where $R'_{n-2}^{\,[n-1]}=(T_{n-3}{+}1)\cdots(T_1{+}1)\in H_q(S_{n-2})$.
$R'_{n-2}^{\,[n-1]}$ preserves $E_{n-1}^+\!\mid_{V_\lambda}$
(since $T_1,\ldots,T_{n-3}$ commute with $T_{n-1}$).  By
Lemma~\ref{lem:Phi-equiv} (one $\Phi_X$ at a time),
\[
   R'_{n-2}^{\,[n-1]}\,\Phi_X(V_{\mu_X})
   \;=\;\Phi_X\big(R'_{n-2}^{(\mu_X)}V_{\mu_X}\big),
   \qquad X\in\{A,B,C\}.
\]
Hence
\[
   R'_{n-1}R'_n V_\lambda
   = (T_{n-2}{+}1)\sum_X\Phi_X\big(R'_{n-2}^{(\mu_X)} V_{\mu_X}\big).
\]

\paragraph{Step 3.}  Apply $R'_{n-2}$.  Write
$R'_{n-2}=(T_{n-3}{+}1)\,S$ where $S=(T_{n-4}{+}1)\cdots(T_1{+}1)\in
H_q(S_{n-3})$.  Then $S$ commutes with $T_{n-2}$ (index gap $\ge 2$),
so $S(T_{n-2}{+}1)=(T_{n-2}{+}1)S$, giving
\[
   R'_{n-2}(T_{n-2}{+}1)
   = (T_{n-3}{+}1)\,S\,(T_{n-2}{+}1)
   = (T_{n-3}{+}1)(T_{n-2}{+}1)\,S.
\]
$S\in H_q(S_{n-3})\subset H_q(S_{n-2})$, so by
Lemma~\ref{lem:Phi-equiv},
$S\,\Phi_X(v)=\Phi_X(R'_{n-3}^{(\mu_X)}\,v)$ for each $X$ (the
operator $S$ acts on $V_{\mu_X}$ as $R'_{n-3}^{(\mu_X)}$, which uses
generators $T_1,\ldots,T_{n-4}$).  Hence
\begin{align*}
   B_{n-2}^{(\lambda)} V_\lambda
   &= R'_{n-2}(T_{n-2}{+}1)\sum_X\Phi_X(R'_{n-2}^{(\mu_X)}V_{\mu_X})\\
   &= (T_{n-3}{+}1)(T_{n-2}{+}1)\,S\,\sum_X\Phi_X(R'_{n-2}^{(\mu_X)}V_{\mu_X})\\
   &= (T_{n-3}{+}1)(T_{n-2}{+}1)\sum_X
      \Phi_X\big(R'_{n-3}^{(\mu_X)}R'_{n-2}^{(\mu_X)}V_{\mu_X}\big)\\
   &= (T_{n-3}{+}1)(T_{n-2}{+}1)\sum_X\Phi_X\big(B_{n-3}^{(\mu_X)} V_{\mu_X}\big),
\end{align*}
which is~\eqref{eq:reduction}.
\end{proof}

\section{The $C$-piece is zero}\label{sec:C-zero}

\begin{lemma}[The $C$-piece vanishes]\label{lem:C-zero}
$B_{n-3}^{(\mu_C)} V_{\mu_C}=0$.
\end{lemma}

\begin{proof}
$\mu_C=(2,1^{n-4})$ has $\ell(\mu_C)=(n-2)-1=n-3$, so
$j(\mu_C)=\ell(\mu_C)+1=n-2$.  By Theorem~\ref{thm:muC},
$B_{n-2}^{(\mu_C)} V_{\mu_C}=R'^{(\mu_C)}_{n-2} V_{\mu_C}\subseteq
E_1^-\!\mid_{V_{\mu_C}}$.

Now $B_{n-3}^{(\mu_C)} V_{\mu_C}=R'^{(\mu_C)}_{n-3}\,B_{n-2}^{(\mu_C)}
V_{\mu_C}$.  The rightmost factor of $R'^{(\mu_C)}_{n-3}$ is
$(T_1{+}1)$, which annihilates $E_1^-\!\mid_{V_{\mu_C}}$:
$(T_1{+}1)v=0$ for $v\in E_1^-$ (since $T_1 v=-v$).  Hence
$R'^{(\mu_C)}_{n-3}\,B_{n-2}^{(\mu_C)} V_{\mu_C}=0$.
\end{proof}

\begin{remark}
Geometrically, the third corner pair $\{c_1,c_2\}=\{(1,3),(2,2)\}$
strips two cells from $\lambda$ without consuming the column-$1$
bottom, leaving a hook $\mu_C$ \emph{shorter by one row} than
$\mu_A,\mu_B$.  The shorter $\mu_C$ has its $j$-formula at $j(\mu_C)=
n-2$ rather than $n-3$, so the index $n-3$ in our reduction sits
\emph{below} the $j$-formula and is annihilated by the leading
$(T_1{+}1)$.  This is the structural reason the third branch
contributes nothing.
\end{remark}

\section{Main theorem}\label{sec:main}

\begin{theorem}[Rank-zero for $\lambda=(3,2,1^{n-5})$]\label{thm:main}
For every $n\ge 8$,
\[
   B_{n-2}^{(\lambda)} V_\lambda\subseteq E_1^-\!\mid_{V_\lambda}.
\]
Consequently, $\Pi^{S_n}\!\mid_{V_\lambda}=0$.
\end{theorem}

\begin{proof}
Combine the reduction (Theorem~\ref{thm:reduction}) with the
$j$-formulas for $\mu_A,\mu_B,\mu_C$:

\emph{$A$-piece.}  $\mu_A=(2,2,1^{n-6})\vdash n-2$ with
$\ell(\mu_A)=n-4$.  By Theorem~\ref{thm:muA} (applied at $m=n-2$,
$\ell+1=n-3$), $B_{n-3}^{(\mu_A)} V_{\mu_A}\subseteq
E_1^-\!\mid_{V_{\mu_A}}$.

\emph{$B$-piece.}  $\mu_B=(3,1^{n-5})\vdash n-2$ with
$\ell(\mu_B)=n-4$.  By Theorem~\ref{thm:muB},
$B_{n-3}^{(\mu_B)} V_{\mu_B}\subseteq E_1^-\!\mid_{V_{\mu_B}}$.

\emph{$C$-piece.}  By Lemma~\ref{lem:C-zero},
$B_{n-3}^{(\mu_C)} V_{\mu_C}=0$.

By Lemma~\ref{lem:Phi-E1}, $\Phi_X(E_1^-\!\mid_{V_{\mu_X}})\subseteq
E_1^-\!\mid_{V_\lambda}$ for $X=A,B$ (the $C$ case is trivial since
the input is $0$).  Hence
\[
   \sum_X\Phi_X(B_{n-3}^{(\mu_X)} V_{\mu_X})\subseteq
   E_1^-\!\mid_{V_\lambda}.
\]

Finally, $(T_{n-3}{+}1)$ and $(T_{n-2}{+}1)$ both commute with $T_1$
(index gap $\ge 2$ since $n-3\ge 5,n-2\ge 6$ for $n\ge 8$).  Hence
each preserves $E_1^-\!\mid_{V_\lambda}$.  Applying both:
$B_{n-2}^{(\lambda)} V_\lambda\subseteq E_1^-\!\mid_{V_\lambda}$.

For the second statement, factor
$\Pi^{S_n}=R'_2 R'_3\cdots R'_{n-3}\cdot B_{n-2}^{(\lambda)}$.
The rightmost factor of $R'_{n-3}$ is $(T_1{+}1)$, which
annihilates $E_1^-$.  Hence $\Pi^{S_n}V_\lambda=0$.
\end{proof}

\begin{remark}[On the dimension]
Theorem~\ref{thm:main} does \emph{not} assert
$\dim B_{n-2}^{(\lambda)} V_\lambda=2$ --- only that the image lies
in $E_1^-$, which is all we need for rank-zero.  The dimension is
empirically $2$ (matching the rank-corner formula
$\#\{\text{length-pres corners}\}=2$), and the proof framework
gives an upper bound of $\dim\Phi_A(B_{n-3}^{(\mu_A)} V_{\mu_A})+
\dim\Phi_B(B_{n-3}^{(\mu_B)} V_{\mu_B})\le 1+1=2$, but lower bound
($\ne 0$ for each piece after $(T_{n-3}{+}1)(T_{n-2}{+}1)$) would
require a Hoefsmit case analysis we leave for a future writeup.
The structural significance is the $E_1^-$ containment.
\end{remark}

\section{Computational verification}\label{sec:verify}

All checks at $q=7$ in $\mathbb{Q}$ (rational arithmetic, no
floating point).

\paragraph{Dimension decomposition.}  For each $n\in\{8,9,10,11\}$,
$\dim V_{\mu_A}+\dim V_{\mu_B}+\dim V_{\mu_C}$ equals
$\dim E_{n-1}^+\!\mid_{V_\lambda}=\dim R'_n V_\lambda$:

\begin{center}
\begin{tabular}{ccccccc}
\toprule
$n$ & $\dim V_\lambda$ & $\dim V_{\mu_A}$ & $\dim V_{\mu_B}$ & $\dim V_{\mu_C}$ & sum & $\dim R'_n V_\lambda$ \\
\midrule
$8$  & $64$  & $9$  & $10$ & $5$ & $24$ & $24$ \\
$9$  & $105$ & $14$ & $15$ & $6$ & $35$ & $35$ \\
$10$ & $160$ & $20$ & $21$ & $7$ & $48$ & $48$ \\
$11$ & $231$ & $27$ & $28$ & $8$ & $63$ & $63$ \\
\bottomrule
\end{tabular}
\end{center}

\paragraph{$E_1^-$ containment at $j=\ell+1$.}  For each
$n\in\{8,9,10\}$, $B_{n-2}^{(\lambda)} V_\lambda$ has dimension $2$
\emph{and} $(T_1{+}1)\cdot B_{n-2}^{(\lambda)} V_\lambda=0$
(equivalent to $B_{n-2}^{(\lambda)} V_\lambda\subseteq E_1^-$):

\begin{center}
\begin{tabular}{ccccc}
\toprule
$n$ & $\dim R'_n V_\lambda$ & $\dim R'_{n-1}R'_n V_\lambda$ & $\dim B_{n-2}V_\lambda$ & in $E_1^-$? \\
\midrule
$8$  & $24$ & $8$  & $2$ & yes \\
$9$  & $35$ & $10$ & $2$ & yes \\
$10$ & $48$ & $12$ & $2$ & yes \\
\bottomrule
\end{tabular}
\end{center}

(At $n=11$ the rank-$8$ step in $R'_{n-1}R'_n V_\lambda$ would land
at $14$; we verified $B_{n-2}V_\lambda$ has dim $2$ in $E_1^-$ via
prior data in \texttt{2026-05-08-conj-81-test/out\_321111\_n11.log}.)

\paragraph{Final collapse.}  At every tested $n$, one more
$R'_{\ell}=R'_{n-3}$ application drops the rank to $0$, confirming
$\Pi^{S_n}\!\mid_{V_\lambda}=0$.

Scripts: \texttt{\textasciitilde/projects/scratch/2026-05-20-321-family/extend\_j.py,
verify\_decomp.py}.

\section{Honest assessment}\label{sec:honest}

\paragraph{What is rigorously proved.}
\begin{enumerate}[leftmargin=2em,topsep=0.3em,itemsep=0.2em]
\item Theorem~\ref{thm:reduction}: the three-branch reduction
formula, conditional on Phase A at $\lambda$
(Theorem~\ref{thm:phaseA}, abstract May-19 result, fully
unconditional).
\item Theorem~\ref{thm:main}: rank-zero theorem for
$\lambda=(3,2,1^{n-5})$ at every $n\ge 8$, unconditional once we
cite Theorems~\ref{thm:muA}--\ref{thm:phaseA} (all rigorously
established in prior writeups).
\end{enumerate}

\paragraph{What is sketched but not proved.}
The dimension claim $\dim B_{n-2}^{(\lambda)} V_\lambda=2$ (matching
the rank-corner formula).  We prove $\le 2$ from the reduction; a
non-vanishing argument for each of the $A$- and $B$-pieces after
$(T_{n-3}{+}1)(T_{n-2}{+}1)$ application would close this.  This is
a finite Hoefsmit case analysis analogous to May-15
Section~5; we do not carry it out here because the $E_1^-$
containment (which is what the rank-zero theorem requires) is
already proved.

\paragraph{Structural significance.}  This is the first rank-zero
proof where $B_{\ell+1}^{(\lambda)} V_\lambda$ is \emph{not}
$1$-dimensional.  The two length-preserving corners of $\lambda$
each contribute an independent $1$-dimensional image in
$E_1^-$, giving a $2$-dimensional total that nonetheless
lies in $E_1^-$.  This generalizes the rank-corner formula
\emph{mechanism}: each length-preserving corner contributes one
``L/R-style'' choice, with the inner $\mu_X$ providing the support
structure.  The structural prediction is that for general
$\lambda$ with $r$ length-preserving corners (rank-corner
formula gives $\dim B_{\ell+1}^{(\lambda)} V_\lambda=r$),
the image decomposes as a sum of $r$ pieces, each arising from
the $j$-formula at a doubly-stripped $\mu_X$.  Pairs of corners
whose ``column-$1$ bottom'' is unavailable (length-pres + length-pres
pairings) contribute $0$ by the same $E_1^-$-leakage mechanism
we used for $\mu_C$ here.

\paragraph{Next steps.}
\begin{enumerate}[leftmargin=2em,topsep=0.3em,itemsep=0.2em,
label=(\arabic*)]
\item Prove the dim-$2$ lower bound for $(3,2,1^{n-5})$ (the missing
piece in Remark above).
\item Test the general $r$-corner conjecture on the next family
with $r\ge 2$: $(4,2,1^{n-6})$ or $(3,3,1^{n-6})$ (both have $2$
length-preserving corners; $(4,2,1^{n-6})$ should give dim $2$,
$(3,3,1^{n-6})$ should give dim $1$ since $\lambda_1=\lambda_2$ so
$r=1$).
\item The rank-corner formula itself --- still conjectural for
general rank-zero $\lambda$ --- gains substantial evidence from
this proof, since it predicts a clean structural answer (number of
branches = rank).
\end{enumerate}

\paragraph{What this closes in the rank-zero programme.}
After May-19 (the $(2^a,1^{n-2a})$ family is unconditional), the
rank-zero theorem holds structurally for:
\begin{itemize}[leftmargin=2em,topsep=0.3em,itemsep=0.2em]
\item All hook families $(k,1^{n-k})$ in rank-zero regime
(May-11, May-12).
\item All $(2^a,1^{n-2a})$ families at $n\ge 3a$ (May-13 through
May-19).
\item Today: $(3,2,1^{n-5})$ at $n\ge 8$ (first \emph{two}-corner family).
\end{itemize}
The next natural targets, in rank-zero regime: $(4,2,1^{n-6})$,
$(3,3,1^{n-6})$, $(4,3,1^{n-7})$, $(3,2,2,1^{n-7})$, and so on.
The three-branch reduction framework here generalizes to each:
identify the length-preserving corners, set up $\Phi_X$ for each
pair $\{c_X^-,c_X^+\}$ including $c_3=$bottom-column-$1$, and apply
the relevant $j$-formula at $\mu_X$.

\end{document}
